
The meta-method selection hypothesis is decomposed into three testable sub-hypotheses, each addressing one causal mechanism component. This staged design isolates failure modes: if data collection succeeds but correlation fails, the issue is feature-method relationships rather than data availability; if correlation succeeds but training fails, the issue is sample size rather than feature informativeness.

\textbf{Hypothesis Decomposition.} The core hypothesis states that a meta-classifier trained on 50-60 benchmarks using fast-to-compute features will predict method families achieving top-30\% ranking with greater than 50\% success rate. This embeds four assumptions: (A1) Published benchmarks aggregate at least 50 training examples. (A2) Fast features (under one minute computation) capture dataset characteristics differentiating method performance. (A3) Random Forest extracts generalizable relationships rather than encoding domain folklore. (A4) Predicted methods achieve competitive performance on held-out datasets.

These are tested sequentially through three sub-hypotheses:

\textbf{H-E1 (Existence).} At least 50 benchmarks with complete baseline comparisons are collectible from literature (OGB, FedML, LEAF, pFL-Bench, Champneys, Zhou suites), validating A1.

\textbf{H-M1 (Mechanism - Correlation).} Dataset characteristics (sample size, dimensionality, class imbalance, signal properties) show significant correlations ($\rho$ > 0.3, p < 0.05) with method family rankings, validating A2 and testing the first causal link.

\textbf{H-M2 (Mechanism - Training).} Meta-classifier trained on aggregated benchmarks achieves at least 30\% accuracy on held-out datasets via leave-5-out cross-validation, validating A3 and A4.

This decomposition is falsifiable at each stage. If h-e1 collects fewer than 30 benchmarks, A1 fails (insufficient training data). If h-m1 finds zero correlations, A2 fails (features non-informative). If h-m2 achieves below 30\% accuracy, A3 or A4 fails (no learning or poor generalization). The staged design prevents misattribution: distinguishing ''not enough data'' from ''data exists but is uninformative'' from ''features are informative but classifier does not learn.''

\textbf{H-E1: Benchmark Data Collection.} The objective is verifying that at least 50 benchmarks with complete method comparisons are accessible from published sources. Six benchmark suites were targeted: OGB (15 graph datasets with GCN, GraphSAINT, ClusterGCN, GraphSAGE families, accessible via Python library), FedML (6 federated learning datasets with FedAvg, FedProx, FedOpt baselines, accessible via GitHub), LEAF (5 datasets with heterogeneity statistics and baseline results, available via GitHub), pFL-Bench (8 personalized federated learning benchmarks, accessible via Papers with Code), Champneys et al. NLSI (5 nonlinear system identification tasks with NARX, LSTM, Polynomial baselines, results reported in paper tables), and Zhou et al. Medical FL (9 medical imaging datasets with FedAvg, SCAFFOLD, FedDyn, augmentation method results, reported in paper).

For each source, the collection protocol involved: (1) API/library access to load datasets via official libraries (OGB, FedML) and extract metadata from loaded objects, (2) GitHub repository cloning or README fetching to document dataset statistics and baseline results, (3) paper table extraction for Champneys and Zhou to manually extract benchmark names, dataset characteristics if reported, and method rankings into structured CSV format. Metadata fields collected per benchmark included Tier 1 features (universal, under one second: sample\_size, dimensionality, num\_classes, class\_imbalance) and Tier 2 features (domain-specific, under one minute: autocorrelation\_lag1, edge\_density, feature\_correlation\_rank), plus method rankings as percentile ranking (0-100\%) for representative methods from each family.

Success criteria: PASS requires at least 50 benchmarks collected with at least 80\% of Tier 1 features populated. PARTIAL indicates 30-49 benchmarks or below 80\% feature completeness. FAIL is fewer than 30 benchmarks (insufficient training data, hypothesis untestable). Collected benchmarks should span domains (vision, NLP, tabular, time-series, graph), scales (small n<1K, medium $1K\leq n$<10K, large $n\geq 10K)$, task types (classification, regression), and structural properties (balanced vs. imbalanced, low-dimensional vs. high-dimensional).

\textbf{H-M1: Feature-Method Correlation Analysis.} The objective is testing whether dataset characteristics correlate with method family performance, validating that fast-to-compute features capture information relevant for method selection. Features are computed in two tiers: Tier 1 (universal, under one second) includes sample\_size (total training samples), dimensionality (number of input features), num\_classes (output classes for classification), class\_imbalance (Gini coefficient of class distribution), and signal\_noise\_ratio (feature variance divided by residual variance for regression). Tier 2 (domain-specific, 5-15 seconds) includes autocorrelation\_lag1 (first-order autocorrelation for time-series), edge\_density (average degree divided by max degree for graphs), and feature\_correlation\_rank (rank of feature correlation matrix for tabular).

Method family representatives map diverse methods to four broad families: Linear (ARX, Ridge Regression, Linear SVM), Polynomial (NARX-Poly, Polynomial Regression, Kernel SVM with polynomial kernel), RNN (LSTM, GRU, vanilla RNN), and Augmentation (data augmentation techniques: DDPM+LS, MixUp, CutMix). For each benchmark, the ranking percentile (0-100\%) of the best-performing method in each family is extracted.

The correlation protocol computes, for each (feature, method\_family) pair: (1) filter benchmarks to those with non-missing values for both feature and method ranking, (2) compute Spearman's rank correlation $\rho$ (robust to outliers and nonlinear monotonic relationships), (3) report p-value with threshold p<0.05 for significance, (4) require |$\rho$|>0.3 for ''meaningful'' correlation. All combinations are tested: 7 features $\times$ 4 method families = 28 tests. Bonferroni correction would require p<0.05/$28\approx 0.0018$, but uncorrected p-values are reported (exploratory analysis).

Expected patterns from literature include: sample\_size $\leftrightarrow$ Augmentation (negative correlation: small datasets benefit more), dimensionality $\leftrightarrow$ Linear (positive: high-dimensional favors regularized linear models), autocorrelation\_lag1 $\leftrightarrow$ RNN (positive: temporal structure favors recurrent models), and class\_imbalance $\leftrightarrow$ Augmentation (positive: imbalanced datasets benefit from synthetic samples).

Success criteria: PASS requires at least 5 significant correlations ($\rho$>0.3, p<0.05) spanning at least 3 method families (evidence that features capture method-relevant information). PARTIAL indicates 3-4 correlations or correlations confined to single family (suggests feature insufficiency). FAIL is fewer than 3 correlations (features non-informative; either need Tier 3 expensive probing or dataset characteristics do not predict method performance).

\textbf{H-M2: Meta-Classifier Training and Evaluation.} The objective is testing whether a Random Forest trained on collected benchmarks can predict method family rankings on held-out datasets with at least 30\% top-30\% success rate. Random Forest (scikit-learn RandomForestClassifier) is used for interpretability and robustness. Input is a 7-dimensional feature vector (Tier 1 + Tier 2 features, domain-conditional). Output is 4-class prediction (Linear, Polynomial, RNN, Augmentation). Hyperparameters include n\_estimators=100, max\_depth=10 (prevent overfitting on small training set), and min\_samples\_split=5 (require at least 5 samples to split node).

The target variable per benchmark is the best-performing method family (highest ranking percentile among the four families), converting regression (predicting exact ranking) to 4-class classification (predicting family). Cross-validation uses leave-5-out CV with 10 folds (if 50 benchmarks collected, each fold trains on 45 and tests on 5), stratified by domain if possible.

Evaluation metrics: Primary is top-30\% success rate (percentage of held-out benchmarks where predicted family achieves at most 30th percentile ranking). Secondary metrics include confusion matrix, per-domain accuracy, and feature importance (SHAP).

Baseline comparisons: Random selection (uniform random choice from four families, expected 30\% top-30\% success by definition), majority class (always predict most frequent winner in training set, expected 30-40\%), and domain folklore (predict based on domain only: $vision\rightarrow RNN$ under pre-Transformer assumption, time-$series\rightarrow RNN, tabular\rightarrow Linear$, expected 40-50\%).

Success criteria: PASS requires at least 35\% top-30\% success rate, significantly better than random (Chi-square test p<0.05). PARTIAL indicates 30-35\% (marginally better than random; suggests weak signal; investigate feature sufficiency or increase training data). FAIL is below 30\% (no better than random) or worse than majority-class baseline (indicates no learning; either features insufficient, sample size too small, or method rankings too chaotic).

\textbf{Mock Data Elimination Protocol.} To ensure reported results use real data rather than hard-coded defaults, verification includes: (1) coverage reporting (fraction of benchmarks with real non-NaN values per feature; 0\% indicates no extraction), (2) variance check (standard deviation per feature; zero variance indicates all values identical, likely mock data artifact), and (3) external review (Phase 4 validation includes manual inspection of feature arrays to verify diversity). This transparency allows readers to assess whether negative results stem from data quality issues versus hypothesis failure.

\textbf{Design Rationale.} The three-stage decomposition (h-e1 $\rightarrow$ h-m1 $\rightarrow$ h-m2) isolates causal chain components. If stage 1 fails, benchmark data is inaccessible (infrastructure problem). If stage 1 passes but stage 2 fails, dataset characteristics do not correlate with method rankings (hypothesis problem). If stages 1-2 pass but stage 3 fails, correlations exist but are not strong enough for prediction (sample size or feature richness problem). This design distinguishes data collection bottleneck (actual finding: h-e1 collected 29 benchmarks with 14\% feature coverage) from weak feature-method relationships (untested due to insufficient data) from meta-classifier learning failure (occurred but attributable to degenerate input).
